\documentclass[11pt]{article}

% Change "review" to "final" to generate the final (sometimes called
% camera-ready) version.
% Change to "preprint" to generate a non-anonymous version with page numbers.
\usepackage[final]{acl}

% Standard package includes
\usepackage{times}
\usepackage{latexsym}

% For proper rendering and hyphenation of words containing Latin
% characters (including in bib files)
\usepackage[T1]{fontenc}
% For Vietnamese characters
% \usepackage[T5]{fontenc}
% See https://www.latex-project.org/help/documentation/encguide.pdf
% for other character sets

% This assumes your files are encoded as UTF8
\usepackage[utf8]{inputenc}

% This is not strictly necessary, and may be commented out,
% but it will improve the layout of the manuscript,
% and will typically save some space.
\usepackage{microtype}

% This is also not strictly necessary, and may be commented out.
% However, it will improve the aesthetics of text in
% the typewriter font.
\usepackage{inconsolata}

%Including images in your LaTeX document requires adding
%additional package(s)
\usepackage{graphicx}

\usepackage{amsmath}
\usepackage{booktabs}

% If the title and author information does not fit in the area
% allocated, uncomment the following
%
%\setlength\titlebox{<dim>}
%
% and set <dim> to something 5cm or larger.

\title{The Role of Swear Words in Toxicity Detection}

% Author information can be set in various styles:
% For several authors from the same institution:
% \author{Author 1 \and ... \and Author n \\
%         Address line \\ ... \\ Address line}
% if the names do not fit well on one line use
%         Author 1 \\ {\bf Author 2} \\ ... \\ {\bf Author n} \\
% For authors from different institutions:
% \author{Author 1 \\ Address line \\  ... \\ Address line
%         \And  ... \And
%         Author n \\ Address line \\ ... \\ Address line}
% To start a separate ``row'' of authors use \AND, as in
% \author{Author 1 \\ Address line \\  ... \\ Address line
%         \AND
%         Author 2 \\ Address line \\ ... \\ Address line \And
%         Author 3 \\ Address line \\ ... \\ Address line}

\author{Eugene Wong \\
  University of Gothenburg \\
  LT2326 Machine Learning for Statistical NLP: Advanced \\
\texttt{guswonju@student.gu.se} \\}

%\author{
%  \textbf{First Author\textsuperscript{1}},
%  \textbf{Second Author\textsuperscript{1,2}},
%  \textbf{Third T. Author\textsuperscript{1}},
%  \textbf{Fourth Author\textsuperscript{1}},
%\\
%  \textbf{Fifth Author\textsuperscript{1,2}},
%  \textbf{Sixth Author\textsuperscript{1}},
%  \textbf{Seventh Author\textsuperscript{1}},
%  \textbf{Eighth Author \textsuperscript{1,2,3,4}},
%\\
%  \textbf{Ninth Author\textsuperscript{1}},
%  \textbf{Tenth Author\textsuperscript{1}},
%  \textbf{Eleventh E. Author\textsuperscript{1,2,3,4,5}},
%  \textbf{Twelfth Author\textsuperscript{1}},
%\\
%  \textbf{Thirteenth Author\textsuperscript{3}},
%  \textbf{Fourteenth F. Author\textsuperscript{2,4}},
%  \textbf{Fifteenth Author\textsuperscript{1}},
%  \textbf{Sixteenth Author\textsuperscript{1}},
%\\
%  \textbf{Seventeenth S. Author\textsuperscript{4,5}},
%  \textbf{Eighteenth Author\textsuperscript{3,4}},
%  \textbf{Nineteenth N. Author\textsuperscript{2,5}},
%  \textbf{Twentieth Author\textsuperscript{1}}
%\\
%\\
%  \textsuperscript{1}Affiliation 1,
%  \textsuperscript{2}Affiliation 2,
%  \textsuperscript{3}Affiliation 3,
%  \textsuperscript{4}Affiliation 4,
%  \textsuperscript{5}Affiliation 5
%\\
%  \small{
%    \textbf{Correspondence:} \href{mailto:email@domain}{email@domain}
%  }
%}

\begin{document}
\maketitle
\begin{abstract}
  \textit{\textbf{Warning:} this paper contains content that may
  be offensive or upsetting.}

  Does swearing equal toxicity? Simple toxicity detection systems
  often treat swear words as an immediate indicator of toxic
  behavior. However, swearing and toxicity do not always align.
  People swear casually when venting or expressing strong emotions. In
  contrast, offensive remarks, serious harassment, and threats often
  use swear-free language.
  In this paper, we evaluate how the presence or absence of swear
  words affects toxicity classifiers across a sliced evaluation
  framework. We find that while a fine-tuned transformer outperforms
  traditional models by a small margin, it still struggles
  significantly to recognize non-toxic comments when they contain
  casual swearing. Our findings show that model architecture cannot overcome
  training dataset bias, highlighting the need for understanding
  conversational intent rather than reacting to vocabulary.
\end{abstract}

\section{Introduction}

Online forums and social media platforms struggle to create a safe
space for users because it is difficult to moderate content
accurately. Automated moderation systems often use keyword blocklists
as an initial filter, while machine learning classifiers often learn
to treat swear words as a shortcut to toxicity
\citep{zhou-etal-2021-challenges}.

However, in everyday conversations, swearing and toxicity are not the
same thing \citep{pamungkas-etal-2020-really}. People casually swear
for humor, excitement, emphasis, or
even frustration without meaning any harm:
\begin{quote}
  \textit{``This article is fucking amazing!''}
\end{quote}

In contrast, an abusive comment or threat can be phrased in polite
language without any swear words:
\begin{quote}
  \textit{``You better watch your back when you walk home tonight.''}
\end{quote}

We can group these failures into two main categories:
\begin{enumerate}
  \item \textbf{False Alarm:} Non-toxic comments containing casual
    swearing are falsely flagged as toxic.
  \item \textbf{Blind Spot:} Swear-free toxic comments are missed.
\end{enumerate}

To investigate how much the moderation systems are impacted, we
formulate our research question:

\begin{quote}
  \noindent\textit{How does the presence or absence of swear words
  affect the performance of toxicity detection systems?}
\end{quote}

We hypothesize that models make the most mistakes when swearing and
toxicity do not align. However, we expect transformers to outperform
traditional bag-of-words models because they can understand words
in the context of the entire sentence.

\section{Methodology}

\subsection{Dataset}

We use the Kaggle Jigsaw Toxic Comment Classification Challenge
dataset\footnote{\url{https://www.kaggle.com/c/jigsaw-toxic-comment-classification-challenge}},
which consists of human-labeled Wikipedia comments. The
dataset is annotated with six labels: \textit{toxic},
\textit{severe\_toxic}, \textit{obscene}, \textit{threat},
\textit{insult}, and \textit{identity\_hate}.

However, for this research, we are only interested in knowing whether
the comment is toxic or not. Therefore, we merge these six labels
into a single binary class \texttt{is\_toxic}. If any of the six
labels is positive, we define it as toxic (\texttt{is\_toxic} = 1),
and non-toxic (\texttt{is\_toxic} = 0) otherwise.

In the test set, a large number of entries were excluded from
evaluation and left unlabeled (marked with $-1$). We remove these
entries, reducing the test set from 153,164 to 63,978. We notice that
both splits are severely imbalanced, with only about 10\% of the
comments labeled as toxic (10.17\% in the training set and 9.76\% in
the test set). Because no validation set is provided, we split the
training data into an 80\% training set (127,656 comments) and a 20\%
validation set (31,915 comments), stratifying across both
\texttt{is\_toxic} and \texttt{has\_swear} to preserve class and swear
proportions across splits.

\subsection{Data Cleanup}

We first apply standard text normalization by unescaping HTML
entities, removing surrounding quotation marks and whitespace, and
collapsing whitespace and linebreaks. Next, we apply three targeted
preprocessing steps:
\begin{enumerate}
  \item \textbf{Wikipedia Markup Removal:} We strip the section
    headers (\texttt{==Header==}) and reply indentation colons
    (\texttt{:}) to make the text cleaner to process.
  \item \textbf{Token Masking:} We replace web URLs with
    \texttt{[URL]} and IP addresses with \texttt{[IP]}. Masking these
    reduces noise.
  \item \textbf{Collapsing Elongated Words:} There is frequent use of
    repeated letters for emphasis (e.g., \textit{``sooooo''}). We
    collapse three or more identical characters down to two (e.g.,
    \textit{``soo''}) to group elongated words together.
\end{enumerate}

\subsection{Swear Word Lexicon}

To determine whether comments contain swear words, we need a lexicon.
For that, we use the dictionary from
\texttt{noswearing.com}\footnote{\url{https://www.noswearing.com}}.

However, we notice that the list contains some terms that are not
offensive by nature. Thus, we apply a whitelist to explicitly remove
these words from our swear word lexicon:

\begin{enumerate}
  \item \textbf{Sexual Orientation:} Words such as \textit{``gay''},
    \textit{``lesbian''}, and \textit{``queer''} are legitimate identity terms.
  \item \textbf{Anatomical Terms:} Words such as \textit{``penis''},
    \textit{``vagina''}, \textit{``testicle''}, and \textit{``anus''}
    are medical and
    biological terms.
\end{enumerate}

After removing the whitelisted words, our final lexicon contains 341
swear words.

To prevent false substring matches (such as matching \textit{``ass''}
  inside \textit{``classic''} or \textit{``hell''} inside
\textit{``hello''}), we compile the lexicon into a regular expression
using word boundaries (\texttt{\textbackslash b}).

We then define a binary class \texttt{has\_swear}, where
\texttt{has\_swear} = 1 if the comment contains at least one swear
word from the lexicon, and \texttt{has\_swear} = 0 otherwise.

\subsection{Evaluation Framework}

Because the dataset is severely imbalanced, standard aggregate
metrics cannot tell the full story
\citep{rottger-etal-2021-hatecheck}. To isolate the effect of swear
words and test our hypothesis, we slice the test set into a
2$\times$2 matrix based on toxicity (\texttt{is\_toxic}) and whether
the comment contains swear words (\texttt{has\_swear}):

\begin{table}[ht]
  \centering
  \caption{Test set comment distribution across toxicity and swear
  word presence ($N = 63,978$).}
  \label{tab:matrix_2x2}
  \resizebox{\columnwidth}{!}{%
    \begin{tabular}{lrrr}
      \toprule
      & \textbf{With Swear} & \textbf{No Swear} & \textbf{Total} \\
      \midrule
      \textbf{Toxic} & 3,234 (5.05\%) & 3,009 (4.70\%) & 6,243 (9.76\%) \\
      \textbf{Non-Toxic} & 2,874 (4.49\%) & 54,861 (85.75\%) & 57,735
      (90.24\%) \\
      \midrule
      \textbf{Total} & 6,108 (9.55\%) & 57,870 (90.45\%) & 63,978 (100\%) \\
      \bottomrule
    \end{tabular}%
  }
\end{table}

Table~\ref{tab:matrix_2x2} reveals two critical challenges:
\begin{itemize}
  \item \textbf{False Alarm Risk:} Among the 6,108 comments containing
    swear words, toxic comments (3,234, or 52.95\%) and non-toxic
    comments (2,874, or 47.05\%) are almost evenly split. Determining
    the comment's toxicity is essentially doing a coin flip.
  \item \textbf{Blind Spot Risk:} Nearly half of all toxic comments
    (48.20\%, or 3,009 out of 6,243) contain no swear words at all.
\end{itemize}

In contrast to the balanced swear-word distribution in the test set,
the training set exhibits severe distribution skew. 78.22\% of all
swearing comments in the training set are toxic (6,865 toxic vs. 1,912
non-toxic). This imbalance creates bias that swearing means toxicity.

To quantify performance, we focus on two key metrics:

\begin{equation}
  \Delta\text{FPR} = \text{FPR}_{\text{with swear}} -
  \text{FPR}_{\text{no swear}}
\end{equation}
$\Delta\text{FPR}$ measures the \textbf{False Alarm Surge}: how much
does casual swearing inflate false alarms? Ideally, a fair model
should achieve $\Delta\text{FPR} = 0$, indicating that casual
swearing never triggers a false accusation.

\begin{equation}
  \Delta\text{Recall} = \text{Recall}_{\text{with swear}} -
  \text{Recall}_{\text{no swear}}
\end{equation}
$\Delta\text{Recall}$ measures the \textbf{Sensitivity Drop}: how
much does detection drop when a comment contains no swear words?
Ideally, a robust model should achieve $\Delta\text{Recall} = 0$,
indicating that it catches toxic comments without depending on swear words.

\subsection{Models Evaluated}

We evaluate five models spanning from simple heuristics to a
pre-trained transformer:

\begin{enumerate}
  \item \textbf{Majority Baseline:} Always predicts the majority
    class (non-toxic). It serves as an accuracy benchmark (90.24\%)
    to demonstrate why accuracy is a flawed metric for an imbalanced dataset.
  \item \textbf{Random Baseline:} Generates predictions randomly
    according to the class distribution (10\% toxic, 90\% non-toxic).
  \item \textbf{Rule-Based Lexicon Baseline:} Predicts that a comment is
    toxic if it contains at least one word from our swear lexicon
    (\texttt{has\_swear} = 1). This acts as a benchmark for lexical bias.
  \item \textbf{TF-IDF + Logistic Regression:} Text is vectorized
    into the top 25,000 unigrams and bigrams using sublinear term
    frequency scaling, followed by a logistic regression classifier
    with balanced class weights to compensate for class imbalance.
  \item \textbf{DeBERTa-v3-base:} A pre-trained transformer that uses
    disentangled self-attention to process words in relation to their
    surrounding context. Analyzing token length across 10,000 sampled
    comments showed that 82.75\% fit within 128 tokens,
    so we set the maximum sequence length to 128. We fine-tune the model
    for 3 epochs with an effective batch size of 32 (batch size 16 with
    gradient accumulation), a learning rate of $2\times 10^{-5}$, and
    mixed precision (\texttt{bf16}), selecting the final model from the
    checkpoint with the highest $F_1$.
\end{enumerate}

\section{Results}

\subsection{Overall Metrics}

Table~\ref{tab:overall_results} presents the overall evaluation metrics
across all five models.

\begin{table*}[t]
  \centering
  \caption{Overall test set performance across baselines and DeBERTa.}
  \label{tab:overall_results}
  \begin{tabular}{lccccc}
    \toprule
    Model & Accuracy & Precision & Recall & $F_1$ & PR-AUC \\
    \midrule
    Majority Baseline & 0.9024 & 0.0000 & 0.0000 & 0.0000 & 0.0976 \\
    Random Baseline & 0.8219 & 0.1000 & 0.1032 & 0.1016 & 0.0978 \\
    Rule-Based Baseline & \textbf{0.9080} & \textbf{0.5295} & 0.5180
    & 0.5237 & N/A \\
    TF-IDF + Logistic Regression & 0.8722 & 0.4267 & 0.9020 & 0.5793 & 0.7474 \\
    DeBERTa & 0.9051 & 0.5074 & \textbf{0.9369} &
    \textbf{0.6583} & \textbf{0.8074} \\
    \bottomrule
  \end{tabular}
\end{table*}

The Majority Baseline achieves a high accuracy of 90.24\% while its
$F_1$ is 0.00\%. This proves that standard accuracy is misleading on
an imbalanced dataset. A similar pattern is shown by the Random
Baseline ($F_1 = 10.16\%$).

The Rule-Based baseline's $F_1$ of 52.37\% (recall of 51.80\% and
precision of 52.95\%) directly reflects the coin-flip analogy we
described in Section~2.4.

As expected, DeBERTa achieves the highest overall performance with an
$F_1$ score of 65.83\% (an absolute gain of 7.90 percentage points over
TF-IDF + Logistic Regression) and a PR-AUC of 80.74\% (substantially
outperforming TF-IDF's 74.74\%).

\begin{table*}[t]
  \centering
  \caption{Sliced evaluation across toxic comments (Recall) and
  non-toxic comments (False Alarm Rate).}
  \label{tab:slice_results}
  \small
  \begin{tabular}{lcccccc}
    \toprule
    & \multicolumn{3}{c}{\textbf{Toxic Comments (Recall)}} &
    \multicolumn{3}{c}{\textbf{Non-Toxic Comments (False Alarm Rate)}} \\
    \cmidrule(lr){2-4} \cmidrule(lr){5-7}
    Model & With Swear & No Swear & $\Delta\text{Recall}$ & With
    Swear & No Swear & $\Delta\text{FPR}$ \\
    \midrule
    Majority Baseline & 0.0000 & 0.0000 & 0.0000 & 0.0000 & 0.0000 & 0.0000 \\
    Random Baseline & 0.1030 & 0.1034 & $-0.0004$ & 0.1023 & 0.1003 & 0.0020 \\
    Rule-Based Baseline & 1.0000 & 0.0000 & 1.0000 & 1.0000 & 0.0000 & 1.0000 \\
    TF-IDF + Logistic Regression & 0.9536 & 0.8465 & 0.1072 & 0.6500
    & 0.1039 & 0.5461 \\
    DeBERTa & \textbf{0.9706} & \textbf{0.9006} & \textbf{0.0700} &
    \textbf{0.5828} & \textbf{0.0730} & \textbf{0.5098} \\
    \bottomrule
  \end{tabular}
\end{table*}

\subsection{Sliced Evaluation and Lexical Bias}

Table~\ref{tab:slice_results} evaluates our five models across the
four slices of our 2$\times$2 matrix.

The Rule-Based baseline perfectly portrays the failure of using
keywords to determine toxicity. It achieves 100\% recall on toxic
comments with swear words, but 0\% recall on toxic comments without
swear words (a 100.00 percentage point sensitivity drop,
$\Delta\text{Recall}$). This means it completely fails to identify
toxic comments when there are no swear words. In addition, it
falsely flags 100\% of casual swearing comments (100.00\% FPR).

Comparing the TF-IDF baseline against DeBERTa highlights both the
strength and the limits of contextual modeling:
\begin{itemize}
  \item \textbf{Where DeBERTa improves (The Blind Spot):} The TF-IDF
    baseline achieves only 84.65\% recall on swear-free toxicity
    (a 10.72 percentage point sensitivity drop, $\Delta\text{Recall}$).
    DeBERTa uses surrounding context to catch more subtle toxicity,
    lifting recall to 90.06\% and shrinking this gap to 7.00
    percentage points ($\Delta\text{Recall}$).
  \item \textbf{Where both models struggle (False Alarms):} TF-IDF
    wrongly flags 65.00\% of non-toxic comments containing casual
    swearing (a false alarm surge of 54.61 percentage points over its
    swear-free rate, $\Delta\text{FPR}$). While DeBERTa's self-attention
    reduces this false alarm rate to 58.28\%, it still exhibits a
    surge of 50.98 percentage points over swear-free comments
    ($\Delta\text{FPR}$).
\end{itemize}

This confirms our hypothesis that while contextual attention helps
improve the performance of toxicity detection, it still cannot
overcome the strong correlation between toxicity and swear words
learned from a biased dataset.

\section{Error Analysis}

To understand why the traditional bag-of-words model fails and where
DeBERTa shines, we inspect the learned feature weights and analyze
error cases from the test predictions.

\subsection{The Bag-of-Words Trap: Feature Weights}

\begin{table}[ht]
  \centering
  \caption{Top learned Logistic Regression coefficients from the
  TF-IDF baseline.}
  \label{tab:lr_weights}
  \small
  \setlength{\tabcolsep}{4pt}
  \begin{tabular}{lr lr}
    \toprule
    \multicolumn{2}{c}{\textbf{Top Toxic Features}} &
    \multicolumn{2}{c}{\textbf{Top Non-Toxic Features}} \\
    \cmidrule(lr){1-2} \cmidrule(lr){3-4}
    Feature & Weight & Feature & Weight \\
    \midrule
    \texttt{fuck} & $+22.70$ & \texttt{thanks} & $-5.88$ \\
    \texttt{fucking} & $+19.75$ & \texttt{talk} & $-5.01$ \\
    \texttt{stupid} & $+18.67$ & \texttt{thank} & $-4.81$ \\
    \texttt{shit} & $+18.03$ & \texttt{thank you} & $-4.79$ \\
    \texttt{idiot} & $+15.86$ & \texttt{please} & $-4.74$ \\
    \bottomrule
  \end{tabular}
\end{table}

Inspecting the top learned coefficients from the TF-IDF with Logistic
Regression model (Table~\ref{tab:lr_weights}) reveals that swear
words have massive positive weights, while the opposite, polite
expressions, carry much smaller negative weights. As a result, a
single swear word can overpower the rest of the sentence and push the
prediction toward the toxic class regardless of the surrounding context.

\subsection{Where Context Helps: DeBERTa's Rescues}

Out of the 2,874 non-toxic comments containing casual swearing, DeBERTa
successfully corrects 466 comments that the TF-IDF baseline falsely
flags. A clear example is shown below:

\begin{quote}
  \small
  \textit{``Well, a station that loses out on the CW... they may be
  shit-out-of-luck, but they won't close, as too much money is at stake...''}
  \par\noindent
  \textbf{Human:} Non-toxic \\
  \textbf{Has Swear:} Yes \\
  \textbf{TF-IDF:} Toxic ($P(\text{toxic}) = 50.10\%$) \\
  \textbf{DeBERTa:} Non-toxic ($P(\text{toxic}) = 3.76\%$)
\end{quote}

TF-IDF flags the comment as toxic because of the isolated word
\textit{``shit''}, while DeBERTa recognizes the harmless idiom
\textit{``shit-out-of-luck''} in context, predicting with high
confidence that it is non-toxic.

Similarly, out of the 3,009 swear-free toxic comments,
DeBERTa catches 320 comments that the TF-IDF baseline misses entirely:

\begin{quote}
  \small
  \textit{``he was the gayest person alive....''}
  \par\noindent
  \textbf{Human:} Toxic \\
  \textbf{Has Swear:} No \\
  \textbf{TF-IDF:} Non-toxic ($P(\text{toxic}) = 34.23\%$) \\
  \textbf{DeBERTa:} Toxic ($P(\text{toxic}) = 97.07\%$)
\end{quote}

This comment does not contain any swear words, so TF-IDF predicts it
as non-toxic. DeBERTa, however, reads the sentence as a whole and
recognizes the targeted insult, correctly classifying it as toxic.

\subsection{Shared Failures}

Despite DeBERTa's contextual capability, there are still cases where
it fails alongside the TF-IDF model. Across the test set, both models
falsely flag 1,402 non-toxic comments containing swear words (48.78\%
of that slice), and both fail to detect 142 swear-free toxic comments
(4.72\% of that slice):

\begin{enumerate}\setlength{\itemsep}{1pt}\setlength{\parskip}{0pt}
  \item \textbf{Emotional Venting (Shared False Alarm):}
    \begin{quote}
      \small
      \textit{``Well honestly I cant be fucking bothered tonight. So
      sleep wee all u editors ans admits . I'm off to bed. CU anaon''}
    \end{quote}
    \small
    \textbf{Human:} Non-toxic \\
    \textbf{Has Swear:} Yes \\
    \textbf{TF-IDF:} Toxic ($P(\text{toxic}) = 99.43\%$) \\
    \textbf{DeBERTa:} Toxic ($P(\text{toxic}) = 99.38\%$)
    \par
    \normalsize
    In this comment, the swear word is used casually to express an
    emotion. However, both models predict with over 99\% certainty
    that the comment is toxic.

  \item \textbf{Polite Threats (Shared Blind Spot):}
    \begin{quote}
      \small
      \textit{``Nice job inserting your left-wing garbage in Jan
        Brewers article. I removed it and if you vandalize the article
      again you will be blocked from editing on Wikipedia.''}
    \end{quote}
    \small
    \textbf{Human:} Toxic \\
    \textbf{Has Swear:} No \\
    \textbf{TF-IDF:} Non-toxic ($P(\text{toxic}) = 47.34\%$) \\
    \textbf{DeBERTa:} Non-toxic ($P(\text{toxic}) = 9.74\%$)
    \par
    \normalsize
    This hostile comment contains no swear words and adopts
    an administrative tone, successfully deceiving both models. DeBERTa
    assigns only a 9.74\% toxicity probability, failing to detect
    the underlying coercive threat.
\end{enumerate}

\section{Conclusion}

\subsection{Takeaways}

In this paper, we evaluated how the presence or absence of swear
words affects the performance of toxicity detection systems using a
2$\times$2 sliced evaluation framework on 63,978~human-labeled
Wikipedia comments. We establish that swearing $\neq$ toxicity.

Traditional bag-of-words models suffer from severe lexical bias. The
presence of swear words causes the false alarm rate to surge from
10.39\% to 65.00\% (a surge of 54.61 percentage points, $\Delta\text{FPR}$).
It also causes a 10.72 percentage point drop in detection recall
($\Delta\text{Recall}$) when toxic comments are swear-free.

DeBERTa, despite its contextual self-attention, still falsely flags
\textbf{58.28\%} of non-toxic comments containing casual swearing.
This shows that even a contextual transformer cannot overcome lexical
bias when the training distribution is heavily skewed. Because 78.22\%
of comments containing swear words in the training set are
labeled as toxic, DeBERTa learns to rely on swear words as a
shortcut. Effective moderation requires understanding conversational
intent, not reacting to vocabulary.

\subsection{Limitations}

We acknowledge several limitations in our experimental setup:

\begin{enumerate}\setlength{\itemsep}{1pt}\setlength{\parskip}{0pt}
  \item \textbf{Lexicon Coverage and Obfuscation:} Our swear word
    lexicon relies on an unofficial dictionary of 341 words and
    cannot be used as a gold standard. In addition, keyword matching
    cannot detect obfuscation and leetspeak (e.g., \textit{``f*ck''}
    or \textit{``sh1t''}). This means that some comments are
    still falsely classified as swear-free.

  \item \textbf{Domain Specificity:} The dataset comes from Wikipedia
    comments, which exhibit different conversational norms, wording,
    and emotional tone compared to other domains like social media
    platforms or forums.

  \item \textbf{Binary Target Consolidation:} Our experiment is
    simplified by merging the six labels into a single
    \texttt{is\_toxic} class. However, re-evaluating the whole
    experiment under the subcategories of toxicity could potentially
    yield new insights.
\end{enumerate}

\subsection{Future Work}

In order to improve the performance of toxicity detection, we should
focus on improving dataset balance and annotation quality rather
than solely increasing model size.

\begin{enumerate}\setlength{\itemsep}{1pt}\setlength{\parskip}{0pt}
  \item \textbf{Counterfactual Data Augmentation:} Training datasets
    could be augmented by inserting casual swear words into non-toxic
    comments without changing the intent and toxicity (e.g.,
      \textit{``This is awesome''} $\to$ \textit{``This is fucking
    awesome''}). This would teach the models that the presence of swear
    words does not imply toxicity.

  \item \textbf{Profanity-Aware Tokenization:} Standard tokenizers
    struggle with obfuscated swearing (e.g., \textit{``f*ck''} or
    \textit{``sh!t''}). We could explore training specialized domain
    tokenizers to better recognize obfuscated swear word variants as
    unified tokens.

  \item \textbf{Fairness-Regularized Objectives:} We need to adjust
    the loss functions so that penalties account for class
    and slice imbalances.
\end{enumerate}

\bibliography{custom}

\end{document}